\chapter{Data flow}\label{ch:flow}

\section{Entry points}
Data enters the system in two ways. In the \emph{form} path the browser posts a \code{RawIntake} to \code{POST /api/design} or \code{POST /api/explore}. In the \emph{chat} path the browser posts messages, an optional pasted attachment and up to three images to \code{POST /api/chat}; images are resized and re-encoded as JPEG in the browser before upload.

\section{Chat request path}
Figure~\ref{fig:pipeline} shows the full path of a chat request from the scope screen to the JSON response, and the steps are as follows. Stage names in the trace are given in parentheses where they exist.
\begin{enumerate}
  \item \textbf{Scope screen} on the user's own words. An out-of-scope feature returns \code{out_of_scope}; an unsafe request returns \code{refused}. No model is called.
  \item \textbf{Injection flags} on the brief and the attachment are written to \code{guardrail_events} and logged; they do not block.
  \item If neither a Gemini key nor a local intake model is available, the result is \code{error}.
  \item \textbf{PII redaction} of the transcript and the attachment (stage \code{Guardrails}).
  \item \textbf{Snapshot transcription} (stage \code{Snapshot}), only when a vision backend is configured.
  \item \textbf{Intake} by the model (stage \code{Brief}): JSON with an intent, extracted fields and, for images, a verbatim transcription.
  \item \textbf{Intent handling.} \code{unsafe_request} returns \code{refused}; \code{out_of_scope} or \code{not_a_beam_brief} returns \code{out_of_scope}, subject to the override for ``details missing'' described in Chapter~\ref{ch:llm}.
  \item \textbf{Grounding} of every extracted number, and the factored-load check.
  \item \textbf{Normalisation and range validation} (stage \code{Guardrails}). The outcome is a validated \code{BeamInput}, or \code{needs_info}, \code{rejected} or \code{out_of_scope}.
  \item \textbf{Calculator, verifier and reviser loop} (stages \code{Calculator}, \code{Verifier}, \code{Reviser}), as in Figure~\ref{fig:loop}.
  \item \textbf{Report data}, pre-formatted strings from the verified result.
  \item \textbf{Narrative} (stage \code{Narrative}), model-written with the number trace, or the template.
  \item \textbf{Section sketch}, a server-side SVG of the cross-section.
  \item \textbf{Audit write} (stage \code{Audit}) of the full \code{RunResult}, then the JSON response.
\end{enumerate}
The form path starts at step 9: it has no scope screen on free text, no intake and no grounding. Its narrative still uses the model when a key is configured.

\begin{figure}[p]
\centering
% Request pipeline for POST /api/chat. Deterministic stages are blue, model stages orange.
{\hyphenpenalty=10000\exhyphenpenalty=10000
\begin{tikzpicture}[x=1cm,y=1cm,scale=0.92,transform shape,
  main/.style={text width=6.4cm,minimum height=0.8cm}]
  \def\dy{1.28}
  \node[neu,main] (n0)  at (0,0)        {\texttt{POST /api/chat}: messages, attachment, images};
  \node[det,main] (n1)  at (0,-1*\dy)   {1. Scope and unsafe-request screen};
  \node[det,main] (n2)  at (0,-2*\dy)   {2. Injection flags (logged, not blocking)};
  \node[det,main] (n3)  at (0,-3*\dy)   {3. Model available? (key or local intake)};
  \node[det,main] (n4)  at (0,-4*\dy)   {4. PII redaction of transcript and attachment};
  \node[llm,main,dashed] (n5) at (0,-5*\dy) {5. Snapshot transcription (only if a vision backend is set)};
  \node[llm,main] (n6)  at (0,-6*\dy)   {6. Intake: model returns JSON (intent, fields)};
  \node[det,main] (n7)  at (0,-7*\dy)   {7. Intent handling};
  \node[det,main] (n8)  at (0,-8*\dy)   {8. Grounding of every extracted number; factored-load check};
  \node[det,main] (n9)  at (0,-9*\dy)   {9. Normalisation, unit conversion, range validation};
  \node[det,main] (n10) at (0,-10*\dy)  {10. Calculator, verifier, reviser loop};
  \node[det,main] (n11) at (0,-11*\dy)  {11. Report data (pre-formatted strings)};
  \node[llm,main] (n12) at (0,-12*\dy)  {12. Narrative by the model; number trace, up to 3 attempts};
  \node[det,main] (n13) at (0,-13*\dy)  {13. Section sketch (SVG)};
  \node[det,main] (n14) at (0,-14*\dy)  {14. Audit write (SQLite)};
  \node[neu,main] (n15) at (0,-15*\dy)  {JSON \texttt{RunResult}, HTTP 200};

  \foreach \a/\b in {0/1,1/2,2/3,3/4,4/5,5/6,6/7,7/8,8/9,9/10,10/11,11/12,12/13,13/14,14/15}
    \draw[flow] (n\a) -- (n\b);

  % early exits
  \node[neu,text width=3.6cm] (x1) at (7.6,-1*\dy)  {\texttt{refused} or \texttt{out\_of\_scope}};
  \node[neu,text width=3.6cm] (x3) at (7.6,-3*\dy)  {\texttt{error} (no model)};
  \node[neu,text width=3.6cm] (x6) at (7.6,-6*\dy)  {\texttt{error} (model unusable)};
  \node[neu,text width=3.6cm] (x7) at (7.6,-7*\dy)  {\texttt{refused} or \texttt{out\_of\_scope}};
  \node[neu,text width=3.6cm] (x9) at (7.6,-9*\dy)  {\texttt{needs\_info}, \texttt{rejected} or \texttt{out\_of\_scope}};
  \node[det,text width=3.6cm] (x12) at (7.6,-12*\dy) {template narrative (fallback)};
  \foreach \s in {1,3,6,7,9}
    \draw[flow] (n\s.east) -- (x\s.west);
  \draw[flow] (n12.east) -- node[lbl,above] {3 rejections} (x12.west);
  \draw[flow] (x12.south) -- ++(0,-0.62) -| (n13.east);
  % records the audit trail of early exits as well
  \node[lbl,text width=3.6cm,align=center] at (7.6,-14*\dy) {every exit is written to the audit log};
\end{tikzpicture}}

\vspace{4pt}
{\small\detkey{} deterministic code \quad \llmkey{} touches the language model \quad \neukey{} entry, exit and response}

\caption{Request pipeline for \code{POST /api/chat}. Blue stages are deterministic code, orange stages touch the language model, and the dashed stage runs only when a vision backend is configured. Grey boxes on the right are early exits, each returned with HTTP 200 and a \code{status} field.}
\label{fig:pipeline}
\end{figure}

\section{Per-stage trace}
Each run stores the timing of its stages as \code{TraceStep} records (stage, origin \code{deterministic} or \code{llm}, milliseconds, status). A stage that runs more than once, such as the calculator in the reviser loop, is summed. The trace is presentation only: it is never hashed, fingerprinted or compared. The model attempts of the run (model, outcome, milliseconds) are stored in \code{llm_attempts}.

\section{Storage and response}
The full \code{RunResult} and the tool versions are stored as JSON in the \code{runs} table (Chapter~\ref{ch:db}). Chat messages are stored redacted; the attachment is not stored, only its length; images are not stored, only their count. Structured form values are stored as submitted. The response is the \code{RunResult} JSON, and the browser renders every number from it. The Markdown report endpoint rebuilds a report from the stored record, not from the model.
